El Trabajo Práctico Nº 1 permitió construir una unidad aritmético-lógica de
ocho bits y separar el cálculo combinacional de la interfaz física de la
\tech{Basys 3}. Los operandos y el código de operación se cargaban entonces con
switches y pulsadores. El resultado y las banderas se observaban en los LED.
Esta organización permitió conservar la \tech{ALU} y modificar por completo la
forma de suministrarle datos.

En este segundo trabajo la computadora envía dos operandos y un código de
operación por el conversor USB--UART de la placa, y recibe como respuesta el
resultado. La transformación parece, en una primera mirada, un reemplazo de
pines. Sin embargo, el enlace UART transmite un bit por vez, no comparte una
señal de reloj con la FPGA y puede detenerse a mitad de una operación. El
circuito ya no solo debe calcular, ahora debe recordar qué recibió, medir intervalos,
administrar capacidad y regresar al estado inicial del protocolo cuando la
entrada es inválida.

Las máquinas de estado finitas organizan ese comportamiento temporal. Una
máquina valida y reconstruye los bytes de entrada; otra produce las tramas de
salida; una tercera interpreta la secuencia A--B--opcode. Entre productor y
consumidor se interponen buffers \tech{FIFO}. En este caso, su función no es aumentar la
velocidad de la UART, sino impedir que diferencias breves de ritmo se conviertan
en pérdidas o sobrescrituras silenciosas de información.
